\emph{Shares.} By Lemma \claimref{dark-matter.projection-every-towers-content-dilutes-ones} a tower's matter acts on this one per unit of the shared three-volume. With R2 the number density in a tower's own space is $1/V_{D}$ in the common cell unit and with R1 the mass per quantum is common, so, as in the proof of \claimref{dark-matter.total-function-thickness}, a five-dimensional tower counted through $\lambda$ cell radii has $\rho_{5} = m\lambda/V_{4}$ and a three-dimensional tower with sheets $\mu$ apart has $\rho_{3} = m/(\mu V_{2})$, against $\rho_{\mathrm{us}} = m/V_{3}$. The shares are $f_{5} = \lambda V_{3}/V_{4} = 8\lambda/3\pi$ and $f_{3} = V_{3}/(\mu V_{2}) = 4/3\mu$, the mirror's is $1$ (Proposition \claimref{dark-matter.population-count-one-share-half-percent}), and with two towers of each odd dimension $T(\lambda,\mu) = 1 + 8/3\mu + 16\lambda/3\pi$. On $\lambda = \mu = c$ this is $T(c)$ of \claimref{dark-matter.total-function-thickness}; the check confirms the agreement at three values of $c$.

(i) \emph{Premise by premise.} R3 is the only premise that names the counted thickness: R2 fixes the density in the tower's own space and a common cell radius, R1 the mass, the sourcing rule the projection, and none of these contains $\lambda$ or $\mu$. F1 to F4, P3, the costed clause and the imports concern algebras and signatures, which the family does not touch. P1 concerns the one frozen point, the one metric and the orientation of the reversal pair $\{2, 6\}$. P2 (\claimref{dark-matter.p2-indifference}) equates the sector labels the descent deposits, along reversal orbits only; the labels are the residues, orientability and the tower charge, which is $\epsilon N$ for the oriented pair and zero for the odd towers, which carry no asymmetry (\claimref{dark-matter.stored-asymmetry-charge-not-number}), all independent of $(\lambda, \mu)$, and the family keeps $\lambda$ and $\mu$ common within each orbit, so P2's equalities hold (the check confirms equal shares and mirrored residues within the three orbits at two members of the family, and lists the labels). So every premise but R3 holds at every $(\lambda, \mu)$. $T$ is affine in $\lambda$ with slope $16/3\pi > 0$ and decreasing in $\mu$; as $\lambda \to 0$ and $\mu \to \infty$ it tends to $1$ and it is unbounded in $\lambda$, so by continuity it takes every value above $1$. The check brackets $T(1,1) \in (5.36, 5.37)$ and $T(2, \tfrac12) \in (9.7, 9.8)$ with $\tfrac{333}{106} < \pi < \tfrac{355}{113}$, and the monotonicity and limits at sample points with $3 < \pi < \tfrac{16}{5}$. Since P2 as scoped reaches labels, which do not vary, it imposes nothing on $(\lambda,\mu)$.

(ii) \emph{Orbits of cells.} The space of a tower is the span of its generators other than the single time; its Gram matrix (computed with \texttt{GA.py} for the six populated representatives) is definite of rank $3$ for this world and the mirror, $2$ for the three-dimensional pair and $4$ for the five-dimensional pair. An isometry preserves dimension, so a unit ball of rank $4$ or $2$ is not carried to one of rank $3$ by any symmetry, while the members of each reversal orbit have equal rank. Hence a licensed rule equates per-cell fillings only within a reversal orbit, and every licensed equality between an odd tower and this one is between fillings of objects common to all towers. \emph{Common objects.} Let $O$ be such a set, with density $n_{O}$ per unit of the shared three-volume, the same for every tower because the objects are shared. A licensed rule fills each object of $O$ with the same number $k$ of quanta in every tower it equates, so an odd tower holds $kn_{O}$ quanta per unit three-volume. If this tower's baryons are the deposited count, its density is also $kn_{O}$ and the share is $1$. Under the charge reading of \claimref{dark-matter.stored-asymmetry-charge-not-number} this tower's baryons are the asymmetric excess, $\epsilon kn_{O}$ with $0 < \epsilon \le 1$ (an excess cannot exceed the count), while an odd tower, with no asymmetry and no annihilation rule, keeps $kn_{O}$, so its share is $1/\epsilon \ge 1$. Equating with a second set $O'$ in place of $O$ gives the same two values, and a rule that equated a filling of $O$ with one of a different set $O'$ would relate objects in different orbits. So the licensed shares are $1$ and $1/\epsilon$. With $3 < \pi < \tfrac{16}{5}$, $8/3\pi \in (\tfrac56, \tfrac89)$, strictly below $1$, hence below every licensed share for every $\epsilon$; the licensed thicknesses solve $8\lambda/3\pi \in \{1, 1/\epsilon\}$, $\lambda = 3\pi/8 \in (\tfrac98, \tfrac65)$ or $3\pi/8\epsilon \ge 3\pi/8$, so never $\lambda = 1$. For the three-dimensional pair $4/3\mu = 1/\epsilon$ at $\mu = 1$ needs $\epsilon = \tfrac34$, which is why the five-dimensional share is the one that decides. With an annihilation stage the odd share would be $s/\epsilon$ for a survival fraction $s$; $s$ is fixed by the odd towers' couplings, and the series has no rule relating those to this tower's (\claimref{dark-matter.no-coupling-rule-odd-towers-follows}), so no licensed rule supplies $s = 8\epsilon/3\pi$. The common objects named in the statement are the rungs of the one wall, at which every tower ascends its own ladder (\claimref{dark-matter.parallel-boundary-mode-spaces-one-wall}), and the light rays of the common count $N$ (\claimref{dark-matter.stored-asymmetry-charge-not-number}); the argument uses only that they are shared.

(iii) By (i) the thickness is undetermined by the premises other than R3, and by (ii) the licensed rules do not reach R3's value; R3 therefore enters the unit-ball relation as an input, and the closed form $(11\pi + 16)/3\pi$ is the value at that input, a fit. The three-dimensional fraction $2f_{3}/T = 8/(3\mu T)$ inherits the same dependence.
